\scanpage{11}
\mspage{4}
plus haut, \deleted{\unsure{pour $g$}} \inserted{savoir $\tau$}, correspond. \deleted{\unsure{aux}} à
\[
\left\{\begin{aligned}
p&=-1\qquad g=l_0^\alpha l_1^\beta\quad(\alpha+\beta=1),\ \alpha,\beta\in\mathbb Z)\\
q&=-1,\ r=1\qquad\text{ou mieux, }\beta=-1
\end{aligned}\right.
\]

\uline{NB} Si on remplace dans les calculs qui précèdent $L_0$ par $L_1$, alors $g$ est remplacé par $\rho(g)$, $q$ et $r$ restant inchangés. On trouve ainsi, en plus du caractère cyclotomique $\chi$, une \deleted{\unsure{deux}} fonction numérique \inserted{$\beta$} sur le groupe
\[
\operatorname{Aut\,ext}_{\mathrm{lac}}(\overset{\pi}{\overset{\Vert}{\hat\pi_{0,3}}},\unsure{\mathfrak S_3|\mathcal S_3})\overset{\mathrm{def}}=\Gamma,
\]
\deleted{\unsure{donnée}} à valeurs dans $\hat{\mathbb Z}$, et une autre \inserted{$\gamma$} à valeurs dans $\pi/\hat{\mathbb Z}$, en faisant opérer $\hat{\mathbb Z}$ sur $\pi$ par une opération $\varphi$
\[
\varphi(\alpha)\cdot g=l_0^\alpha gl_1^{-\alpha}
\]
L'application
\[
(\gamma,\chi):\quad\Gamma\longrightarrow\pi/\hat{\mathbb Z}\times\hat{\mathbb Z}^{*}
\]
\[
u\longmapsto(\gamma(u),\chi(u))
\]
\margin{\uline{NB} $\gamma_0:\mathcal E_0\longrightarrow\pi$ est un $1$-cocycle sur $\mathcal E_0\unsure{=|\simeq}\operatorname{Aut}_{\mathrm{lac}}(\pi,L_0)$, \deleted{\unsure{groupe}} \inserted{\unsure{qui définit}} le ss-gr. des éléments qui commutent ext. à $\unsure{\mathfrak S_3|\mathcal S_3}$}
est injective, \deleted{\unsure{et}} $\beta$ étant défini par $\gamma$ car $\gamma$ prend ses valeurs dans les classes des $g\in\pi$ (mod $\varphi(\hat{\mathbb Z})$) satisfaisant une relation
\mspage{10}
